% Transcription — fonds Alexandre Grothendieck, shelfmark 150, batch 5 (pages 81-98).
% Established from the facsimile archives/batches/150/batch-05.pdf (a copy of
% 150.pdf fetched through the project's relay,
% https://grothendieck-relay.onrender.com/source/150.pdf; PDF page k =
% archivists' page 80+k, no cover sheet), per the skill transcribe-grothendieck.
% The folder is 98 pages in five batches (1-20, 21-40, 41-60, 61-80, 81-98);
% this is the fifth and last (18 pages), done in parallel with batch 4 (not
% yet written when this pass finished) and aligned on batches 1-3.
%
% Pass: Opus 5.5 (claude-opus-5-5), 2026-09-26 — first pass, unchecked against the pages by a human.
%
% Pages skipped: none. Every page carries his mathematics.
%
% Pages summarised: none. The scribbled-over block at the foot of page 98 is
% given as one \struck{\ill{}} with a \note.
%
% His pagination and numbered runs. The batch continues the run headed
% « Espaces stratifiés (étude par voisinages coniques) »: page 80 (batch 4)
% is his sheet 36, ending on (66) J ≺ L. Sheets numbered top left here:
% 37 (p. 81), 38 (83), 41 (88), 42 (90), 43 (92), 44 (94), 45 (96), 46 (98);
% the unnumbered pages 82, 87, 89, 91, 93, 95, 97 are continuations (probably
% versos). Page 84 carries a « 16 » (underlined) top left, which is the number
% of paragraph 16 « Sigma_J et Sigma_{J,L} », not a sheet number; pages 84-86
% (paragraph 16) carry sheet numbers 39 (p. 85) and 40 (p. 86), so pages 84-87
% interrupt the sheet run between 38 and 41 — page 84 presumably stands for
% the missing number. Paragraph 17 « Récapitulation sur les V_i, ... et
% introduction des V_{J_0,...,J_n} » opens on page 88 and restarts the theory
% of bands and of the domination J ≺ K (84)-(87). Pages 92 and 93 were
% scanned in the wrong order: page 91 (end of sheet 42) is continued by
% page 93 (formulas (92), (93)), and only then by sheet 43 (page 92, (94)-
% (98)); said in \note{}s at pp. 91-93. Numbered formulas (67)-(113), with
% repeats kept as written and noted: (71)-(73) twice (pp. 83, 84), (94)-(98)
% on pp. 92 and again on pp. 94-95; (69) and (84) occur only in marginal
% notes. Conditions (Strat 1)-(Strat 5), (Strat 4 bis) on pp. 90-91. No date
% on these pages.
%
% What the batch is about. Pages 81-83: a lemma characterising pairs of bands
% J ≺ L of an ordered set I (the band I' = {i | j <= i <= l}, (67)); the
% resulting description (68)-(73) of V_J^L as a union / colimit of the
% V^{X'}_d over the flags d in Drap_J^L(I) (first term in J, last in L), with
% the spectral sequence (73). Pages 84-87, paragraph 16: Sigma_J for a band J
% (X*_J minus the open tubes T_i, i in bar J \ J), Sigma_J = lim over
% Drap(J) of Sigma_d (72), the boundary dot Sigma_J (74), and Sigma_{J,L}
% (77)-(81) as the boundary of the tube of Sigma_J in Sigma_{I'} cut by
% Sigma_L, with the maps Sigma_{J,L} -> Sigma_L, Sigma_{J,L} -> Sigma_J; a
% remark that Sigma_{J,L} -> Sigma_L makes sense in tame topology without the
% local conical retractions, while Sigma_{J,L} -> Sigma_J needs them. Pages
% 88-98, paragraph 17: the lower topology on I (closed = initial segments),
% bands as locally closed parts, domination; « ordinary » stratifications of
% type I (Strat 1-3), the semi-simplicial resolution X_d (88)-(89), open
% strata X*_i, fidelity (Strat 4, 4 bis) and density (Strat 5), X*_J and
% X_J (93)-(98); then the pro-tubular neighbourhoods V_i = V(X*_i, X), the
% criterion V_i ∩ X*_j ≠ ∅ iff i <= j, the V_{i_0...i_n} ((98), p. 95) as
% building blocks, V^Z, V_J, V_{J,L}, V_J^L as colimits over sets of flags
% (99)-(111), and V_{J_0,...,J_n} for a chain of bands (113), where the
% folder breaks off.
%
% Notation, following batches 1-3: \mathcal V for his curly V (pro-
% neighbourhoods), Sigma_J, Sigma_d, \mathrm{Drap}(I), \varinjlim for his
% underlined « lim » with an arrow, \overline J (lower closure) and
% \widetilde J (upper closure) as he writes them, X^{*}_J, \dot X, \dot T,
% \overset{\circ}{T}, \mathfrak P(I); (h) over an arrow is his mark for a
% homotopy equivalence, set \overset{(h)}{\hookrightarrow}. « cônique » keeps
% his circumflex where he writes it (p. 87 has both « coniques » and
% « côniques »). His underlinings in running text are \emph{}.
%
% Hardest pages: 82 and 85 (long sideways notes in the left margin, read only
% in fragments), 84 (a crosswise note written over the paragraph title), 87
% and 88 (struck interlinear substitutions), 95 (fast connective prose). The
% formulas and numbered statements carry; much of the sideways marginalia does
% not and is left \ill{}.

\documentclass[11pt,a4paper]{article}
% Le préambule commun à toutes les transcriptions du fonds Grothendieck.
%
% Il tient en une idée : **l'incertitude se compose, elle ne se résout pas.**
% Une transcription qui lisse un mot douteux a détruit la seule chose pour
% laquelle on la faisait — un lecteur doit pouvoir savoir, à chaque endroit,
% ce qui était sur le feuillet et ce qui a été ajouté. Les six macros qui
% suivent sont donc l'appareil critique, et non de la décoration.
%
% Les mêmes commandes sont reconnues par `scripts/render.mjs`, qui produit la
% vue de lecture du site. Une macro ajoutée ici doit l'être aussi là-bas,
% faute de quoi le rendu échouera bruyamment — ce qui est voulu : un rendu
% silencieusement faux est pire qu'un rendu absent.

\ProvidesPackage{grothendieck}[2026/08/08 Transcriptions du fonds Grothendieck]

\usepackage{amsmath,amssymb,amsthm}
\usepackage{mathrsfs}
\usepackage{stmaryrd}
% Les diagrammes commutatifs sont partout dans ce fonds : c'est la langue
% même de la géométrie algébrique de Grothendieck, et une transcription qui
% les rendrait en prose ne transcrirait rien.
\usepackage{tikz-cd}
% Français par défaut — c'est la langue des feuillets — mais l'anglais est
% chargé aussi : la traduction et le résumé sont des documents anglais, et la
% typographie française (espaces avant les deux-points) y serait une faute.
% Ces documents basculent par \selectlanguage{english} en tête de corps.
\usepackage[english,french]{babel}
\usepackage{fontspec}
\usepackage{geometry}
\geometry{a4paper,margin=25mm,marginparwidth=28mm}
\usepackage{marginnote}
\usepackage{xcolor}
% ulem, not soul. `soul`'s \st and \ul are built on a character-by-character
% scan that XeLaTeX's font machinery breaks: the document fails with an
% undefined control sequence rather than degrading. ulem does the same two jobs
% and is Unicode-engine safe. `normalem` stops it hijacking \emph, which the
% transcriptions use for Grothendieck's own emphasis.
\usepackage[normalem]{ulem}
\usepackage{hyperref}
\hypersetup{colorlinks=true,linkcolor=black,urlcolor=black,pdfborder={0 0 0}}

% --- Métadonnées du lot -----------------------------------------------------
% Reprises telles quelles par le script de rendu, qui les affiche en tête de
% la vue de lecture. Elles fixent ce que le fichier prétend couvrir : sans
% elles, une transcription flotte sans point d'attache dans un fonds de seize
% mille feuillets.

\newcommand{\folder}[1]{\gdef\grothFolder{#1}}
\newcommand{\batch}[1]{\gdef\grothBatch{#1}}
\newcommand{\leaves}[2]{\gdef\grothFirst{#1}\gdef\grothLast{#2}}
\newcommand{\foldertitle}[1]{\gdef\grothFoldertitle{#1}}
\gdef\grothFolder{?}\gdef\grothBatch{?}\gdef\grothFirst{?}\gdef\grothLast{?}\gdef\grothFoldertitle{}

% --- Le résumé --------------------------------------------------------------
%
% Il ouvre la lecture modernisée, sous le titre « Résumé », et il est composé
% en petit. Ce n'est pas un ornement : c'est le seul passage du document qui
% ne soit pas des mathématiques, et un lecteur doit pouvoir le reconnaître
% avant de l'avoir lu — puis le sauter s'il connaît déjà le sujet, ou s'y
% arrêter s'il ne le connaît pas. Le filet qui le suit marque la reprise du
% corps.

\newenvironment{resume}
  {\small\setlength{\parskip}{0.45em}}
  {\par\medskip\hrule\medskip}

% --- Mots-clés ---------------------------------------------------------------
%
% En anglais, en clôture du résumé de la lecture modernisée : le vocabulaire
% moderne sous lequel chercher ce que les feuillets construisent. Le manifeste
% du site (`npm run manifest`) les extrait pour étiqueter la cote entière —
% cette ligne est la source unique des tags, il n'y a pas de fichier à côté.
\newcommand{\keywords}[1]{\par\smallskip\noindent
  {\footnotesize\textsc{Keywords} --- \textit{#1}}\par}

% --- L'appareil critique ----------------------------------------------------

% Le numéro de feuillet, en marge. C'est l'ancre qui permet de régler un
% désaccord sur une formule en regardant *un* feuillet plutôt qu'une chemise
% de six cents. Le site s'en sert aussi pour tourner les pages du fac-similé
% quand on fait défiler la transcription.
\newcommand{\leaf}[1]{%
  \marginnote{\footnotesize\color{gray!70}\textsf{#1}}%
  \label{leaf:#1}\ignorespaces}

% Illisible : on ne devine pas. Un mot inventé qui se lit comme les autres est
% le pire résultat possible.
\newcommand{\ill}{\textcolor{red!60!black}{[\,\ldots\,]}}

% Lecture douteuse : le mot est donné, et signalé comme incertain.
\newcommand{\uncertain}[1]{\uline{#1}}

% Ajout de l'éditeur — une lettre manquante, un mot sous-entendu.
\newcommand{\add}[1]{\textcolor{blue!50!black}{[#1]}}

% Biffé par Grothendieck lui-même. Ce qu'il a rayé fait partie du document :
% une rature dit souvent où la pensée a changé de direction.
\newcommand{\struck}[1]{\sout{#1}}

% Note du transcripteur, sur l'état matériel du feuillet.
\newcommand{\note}[1]{\par\smallskip{\small\color{gray!60!black}\textit{[#1]}}\par\smallskip}

% Note marginale de Grothendieck, distincte du corps du texte.
\newcommand{\marginal}[1]{\par\smallskip\noindent\hspace{1em}%
  {\small\color{gray!60!black}#1}\par\smallskip}

% --- Titre ------------------------------------------------------------------

\AtBeginDocument{%
  \begin{center}
    {\large\bfseries Fonds Alexandre Grothendieck --- cote n\textsuperscript{o}~\grothFolder}\\[2pt]
    {\itshape\grothFoldertitle}\\[4pt]
    {\small feuillets \grothFirst--\grothLast\ (lot \grothBatch)}\\[2pt]
    {\footnotesize\color{gray!60!black}Transcription établie d'après le fac-similé numérisé
     par l'Université de Montpellier.\\
     Les lectures incertaines sont soulignées, les illisibles notées [\,\ldots\,],
     les ajouts entre crochets.}
  \end{center}
  \vspace{1em}}

\endinput


\folder{150}
\batch{5}
\pages{81}{98}
\dating{1981-1982}
\watermark{Édition de démonstration}
\foldertitle{Espaces stratifiés et voisinages côniques (ou : déploiement des espaces stratifiés) : notes manuscrites (s.d.)}

\begin{document}

\page{81}
\note{En haut à gauche, « 37 » de sa main : trente-septième feuillet de la suite « Espaces stratifiés (étude par voisinages coniques) » (feuillets 6 à 16 au lot 2, 17 à 26 au lot 3).}

\emph{Lemme}. Soient $I$ un ens.\ ordonné, $J,L$ deux parties de $I$. Conditions équivalentes

a) 1°) $J,L$ sont des bandes, 2°) tout élément de $J$ est majoré par un élément de $L$, tout élément de $L$ majore un élément de $J$, \struck{3°) si $j\in J$, $l\in L$ sont tels que $l\le j$, alors $j,l\in J\cap L$}

\marginal{i.e.\ $J\prec L$}\note{dans la marge gauche, en face de a).}

b) $\exists$ une bande $I'$ de $I$, telle que $J$ soit un idéal de $I'$, $L$ un coidéal de $I'$, et que tout élément de $I'$ majore un élément de $J$ et soit majoré par un élément de $L$.\note{« idéal de $I'$ » : l'accent est écrit avant le $I$.}

Sous ces conditions, \struck{il existe \ill{} plus} \add{la bande $I'$ de b) est unique, c'est} la plus petite bande $I'$ contenant $J$ et $L$ satisfait b), et est donnée par
\[
(67)\qquad I'=\lbrace i\in I \mid \exists j\in J,\ l\in L \text{ avec } j\le i\le l\rbrace
\]
\note{le numéro, écrit en marge, se lit mal ; « (67) » est suggéré par le (68) de la page suivante.}

\emph{Dém.} L'implication b) $\Rightarrow$ a) triviale, (prouvons \add{ainsi que le fait que $I'$ est donnée par (67)}) a) $\Rightarrow$ b). Soit $I'$ donné par (67), on a $I'\supset J\cup L$ (par a) 2°)) $I'$ est une bande (trivial), et c'est évidemment la plus petite bande $\supset J\cup L$. Prouvons que $J$ est un idéal de $I'$ i.e.\ $j\in J$, $i\in I'$, $i\le j\Rightarrow i\in J$. Comme $i\in I'$, $\exists j_0\in J$ tel que $j_0\le i$, donc $j_0\le i\le j$ avec $j_0,j\in J$, et comme $J$ est une bande \add{(par a) 1°)}, $i\in J$. \struck{\ill{}} On voit de même que $L$ est un coidéal de $I'$. Soit $i\in I'$, \struck{\ill{}} \add{par définition} de $I'$ $i$ majore un élément de \uncertain{$I$} et est majoré par un élément de $L$, cqfd.\note{on attend « de $J$ ».}

\page{82}
\note{Sans numéro de sa main ; suite du feuillet 37.}

Dans le cas actuel, on trouve donc

(68) $\mathcal V_J^L\simeq$ le \ill{} de la forme $\mathcal V'^{\,L}_J$, relatif à l'espace $X'=X^{*}_{I'}$, stratifié par les $X'_i=X_i\cap X^{*}_{I'}$ ($i\in I'$), où $X'_J$ est fermé, et $X'_L$ est ouvert, \struck{donc}\note{le mot après « le » commence par une lettre biffée ; il n'a pas pu être lu. Le signe devant « le » est un $=$ surmonté d'un trait ondulé.}

On trouve
\[
\mathcal V_J^L=\mathcal V_J\cap X^{*}_L=\Big(\bigcup_{j\in J}\mathcal V_j^{X'}\Big)\cap\underbrace{\Big(\bigcup_{l\in L}\mathcal V_l^{X'}\Big)}_{=X^{*}_L}=\bigcup_{j\in J,\ l\in L}\mathcal V^{X'}_{j,l}
\]
\note{l'accolade est tracée au-dessus de la seconde parenthèse, avec l'annotation « $X^{*}_L$ car $X^{*}_L$ ouvert dans $X'$ ».}

Notons que
\[
\mathcal V^{X'}_{j,l}=
\begin{cases}
\emptyset & \text{si } j,l \text{ non comparables}^{*}\\
\mathcal V^{X'}_j & \text{si } j=l, \text{ cas dans } j,l\in J\cap L\\
\mathcal V^{X'}_{j,l} \text{ standard} & \text{si } j<l\\
\mathcal V^{X'}_{l,j} \text{ standard} & \text{si } l<j \text{ dans } l,j\in J\cap L
\end{cases}
\]
(dans $l\in J$, $j\in L$\,!)\note{« cas dans » et « dans » : lecture incertaine ; la dernière parenthèse est écrite sous la dernière ligne de l'accolade.}

\marginal{${}^{*}$ Pour raison \ill{} les ouverts de $X$ étant \ill{} \ill{}, \ill{} $X$ \uncertain{unibranche}\ldots}\note{note de renvoi écrite en travers dans la marge gauche ; lecture très incertaine.}

\[
\mathcal V_J^L=\bigcup_{\substack{j\in J\\ l\in L\\ j\le l}}\mathcal V^{X'}_{j,l}
\]
pour \ill{}
(où on a appelé que $\mathcal V^{X'}_{j,j}=\mathcal V^{X'}_j$)

\marginal{NB Pour cette formule, \ill{} \ill{}, valable si $J\prec L$, on a (69) \ill{} (équiv.\ d'homotopie) \ldots\ $\mathcal V_J^L\xrightarrow{(h)}\mathcal V_J\cap\mathcal V_L\overset{\text{déf}}{=}\mathcal V_{J,L}$, où $\mathcal V_J^L\overset{\text{déf}}{=}\mathcal V_J\cap X^{*}_L$}\note{longue note en travers dans la marge gauche, en face de la formule ; lecture très incertaine, sauf les formules. Le numéro (69) figure dans cette note.}

Pour écrire $\mathcal V_J^L$ par \emph{recollement} des $\mathcal V^{X'}_{j,l}$, il faut expliciter les
\[
\mathcal V^{X'}_{j,l}\cap\mathcal V^{X'}_{j',l'}
\]
qui sont vides, sauf si $\lbrace j,j',l,l'\rbrace$ est un drapeau. C'est un drapeau \struck{dont} ayant au plus quatre termes \add{(distincts)}, ceux-ci \struck{qui} sont dans $J\cup L$, et il en est au \struck{deux dans $J$, deux dans $L$, et au moins, un} moins un dans $J$, au moins un dans $L$. \uncertain{Ce sont} donc des drapeaux de la forme

\page{83}
\note{En haut à gauche, « 38 » de sa main.}

\[
(70)\quad
\begin{cases}
\mathcal V^{X'}_i & i\in J\cap L\\
\mathcal V^{X'}_{ij} & i\in J\smallsetminus J\cap L,\ j\in L\smallsetminus J\cap L \qquad (i<j)\\
\mathcal V^{X'}_{ijk} & i\in J,\ k\in L,\ j\in J\cup L \qquad (i<j<k)\\
\mathcal V^{X'}_{ijkl} & i,j\in J,\ k,l\in L \qquad (i<j<k<l)
\end{cases}
\]
\note{dans la deuxième ligne, « $\smallsetminus J\cap L$ » après $i\in J$ est ajouté au-dessus de la ligne.}

\marginal{NB le cas le plus intéressant est celui où $J\cap L=\emptyset$ \ldots\ \ill{} \ill{} \ill{} \ill{} \ill{} $\mathcal V^{X'}_i$ \ill{} \ldots\ il \ill{} \ill{} \ill{} \ill{} cas \ill{}}\note{notes en travers dans la marge gauche, en face de (70) ; lecture très incertaine.}

On peut donc dire que $\mathcal V^L_J$ est la $\varinjlim$ (réunion) des prolongements précédents, où \ill{} \struck{\ill{}} \add{entre ces espaces} les morphismes semi-simpliciaux \struck{qui en explicite}.

\struck{Ceci \ill{} \ill{} \ill{}} Si on \ill{} \add{les $\mathcal V_{ijk}$ stables par intersections} dans un système \struck{semi-simplicial} d'espaces, qui par \struck{recollement} $\varinjlim$ donne $\mathcal V^L_J$, on prendra les
\[
(71)\qquad \mathcal V^{X'}_{i_0\ldots i_n}\quad\text{avec } i_0<\cdots<i_n,\quad
\begin{cases} i_0,\ldots,i_n\in J\cup L\\ i_0\in J,\ i_n\in L\end{cases}
\]
\struck{\ill{} \ill{} \ill{}} les flèches entre ces espaces correspondant aux inclusions de drapeaux \emph{satisfaisant} les \emph{conditions précédentes}. Posons donc
\[
(72)\qquad \mathrm{Drap}^L_J(I)\subset\mathrm{Drap}(J\cup L)\subset\mathrm{Drap}(I)
\]
« drapeaux $(i_0\ldots i_n)$ tels que $i_0\in J$, $i_n\in L$ ».\note{le $J$ de $\mathrm{Drap}(J\cup L)$ est récrit par-dessus une autre lettre.}

On a donc un recouvrement de $\mathcal V^L_J$ par les $\mathcal V^{X'}_d$, indexé par l'ens.\ ordonné précédent. On aura p.\ ex., pour un faisceau $F$ sur $\mathcal V^L_J$
\[
(73)\qquad H^{*}(\mathcal V^L_J,F)\Longleftarrow E_2^{pq}=H^p\big(\mathrm{Drap}^L_J(I),\ d\mapsto H^q(\mathcal V^{X'}_d,F)\big)
\]

\page{84}
\note{En haut à gauche, « 16 » souligné, de sa main : c'est ici un numéro de paragraphe (le paragraphe 17 s'ouvre page 88), non un numéro de feuillet ; la page ne porte pas de numéro de feuillet. Le titre est en partie recouvert par la longue note écrite en travers de la marge.}

\subsection*{16. $\Sigma_J$ et $\Sigma_{J,L}$}

\marginal{Un point de vue plus \ill{} \ill{} \ill{} \ill{} les $\Sigma_i$, $\Sigma_{i_0\ldots i_n}$ est \ill{} \ill{} $\Sigma_J$, $\Sigma_J^L$ \ill{} \ill{} \ill{} (70) \ill{} \ill{} \ill{} $X=\varinjlim_{d\in\mathrm{Drap}(I)}\Sigma_d$ \ldots\ $\Sigma_i\supset X^{*}_i$ \ldots\ $\Sigma_J\supset X^{*}_J=\Sigma^{*}_J$ \ldots}\note{longue note écrite en travers du coin supérieur gauche, par-dessus le titre ; seuls quelques fragments se lisent, dans une lecture très incertaine.}

Soit $J$ une bande de $I$, on définit
\[
\Sigma_J\subset X^{*}_J,\qquad \Sigma_J=X^{*}_J\smallsetminus X^{*}_J\cap\bigcup_{i\in\overline J\smallsetminus J}\overset{\circ}{T}_i
\]
\struck{NB $\bigcup_{j\in J}\Sigma_j$, $X^{*}_J\smallsetminus X^{*}_J\cap\bigcup_{i\in\overline J\smallsetminus J}T_i$}\note{ligne barrée de plusieurs traits.}

\marginal{NB $\overline J=$ idéal engendré par $J$ $=\lbrace i\in I\mid \exists j\in J \text{ avec } i\le j\rbrace$}\note{à droite de la définition.}

\marginal{NB On a \ill{} $X_J=\bigcup_{j\in J}X_j$ $\big(=\bigcup_{j\in\overline J}X^{*}_j=X^{*}_{\overline J}\big)$}\note{dans la marge gauche, en travers.}

On a donc
\[
(71)\quad
\begin{cases}
\Sigma_\emptyset=\emptyset,\quad \Sigma_{\lbrace j\rbrace}=\Sigma_j,\quad \Sigma_I=X\\
\Sigma_J=X_J\quad\text{si } J \text{ un idéal i.e. } J=\overline J
\end{cases}
\]
\note{les numéros (71) à (73) de cette page répètent ceux de la page 83 pour d'autres formules ; ils sont gardés tels qu'écrits.}

\ill{}, \uncertain{sauf erreur}, on aura
\[
(72)\qquad \Sigma_J\simeq\varinjlim_{i,j\in J}\big(\Sigma_j,\ \Sigma_{ij},\ (s_{ij},b_{ij}:\Sigma_{ij}\rightrightarrows\Sigma_i,\Sigma_j)\big)\simeq\varinjlim_{d\in\mathrm{Drap}(J)}\Sigma_d
\]
\struck{\ill{} \ill{} \ill{}}\note{dans (72), $s_{ij}$ va de $\Sigma_{ij}$ dans $\Sigma_i$ et $b_{ij}$ dans $\Sigma_j$ ; les deux flèches sont écrites l'une sous l'autre.}

\marginal{NB $\Sigma_J$ \ill{} \ill{} \ill{} (72) \ill{} \ill{} \ill{} $(\Sigma_j,\Sigma_{jk},s_{jk},b_{jk})$ \ldots\ $j,k\in I$}\note{en travers dans la marge gauche ; lecture très incertaine.}

\[
(73)\qquad \Sigma_J\overset{h}{\underset{\sim}{\hookrightarrow}}X^{*}_J
\]
(équivalence d'homotopie)

On pose
\[
(74)\qquad \dot\Sigma_J=X_J\cap\big(\text{bord de }\underbrace{T_{\overline J\smallsetminus J}}_{\bigcup_{i\in\overline J\smallsetminus J}T_i}\big)
=\bigcup_{\substack{i<j\\ i,j\in J}}\mathrm{Im}(\Sigma_{ij}\to\Sigma_j)
=\varinjlim_{\substack{d\in\mathrm{Drap}\\ d \text{ ayant au moins 2 termes}}}\Sigma_d
\]

\page{85}
\note{En haut à gauche, « 39 » de sa main.}

Soient $J$ et $L$ deux bandes \emph{disjointes}, avec $J\prec L$, une bande $I'$ telle que $J$ soit un idéal de $I'$, $L$ un coidéal de $I'$, $I'\subset\widetilde J\cap\overline L$ (en fait \ill{} \ill{} \ill{} $I'=\widetilde J\cap\overline L$). Soit $J'=I'\smallsetminus L$, donc dans $I'$ on a une suite d'idéaux
\[
J\subset J'\subset I'
\]
(car $J\cap L=\emptyset$). Considérons \struck{$X^{*}_J\subset X^{*}_{J'}$} le diagramme

\begin{tikzcd}[column sep=small]
X_J=X^{*}_{\overline J} \arrow[r, hook] & X_{J'}=X^{*}_{\overline{J'}} \arrow[r, hook] & X_{I'}=X^{*}_{\overline{I'}} \\
X_{J_0}=X^{*}_{J_0} \arrow[u, hook] \arrow[r, no head] & X_{J'_0}=X^{*}_{J'_0} \arrow[u, hook] \arrow[r, hook] & X_{I'_0}=X^{*}_{I'_0} \arrow[u, hook]
\end{tikzcd}

correspondant au diagramme

\begin{tikzcd}[column sep=small]
\overline J \arrow[r, hook] & \overline{J'} \arrow[r, hook] & \overline{I'} \\
J_0 \arrow[u, hook] \arrow[r, hook] & J'_0 \arrow[u, hook] \arrow[r, hook] & I'_0 \arrow[u, hook]
\end{tikzcd}

où
\[
J_0=\overline J\smallsetminus J,\quad J'_0=\overline{J'}\smallsetminus J',\quad I'_0=\overline{I'}\smallsetminus I'.
\]
\note{dans les deux diagrammes, les deux carrés portent la mention « cart » ; ses inclusions sont tracées comme un trait à crochet sans pointe, et les traits verticaux, qui se terminent en bas par un crochet, sont lus ici comme les inclusions $X_{J_0}\hookrightarrow X_J$, etc. ; le premier trait horizontal du bas est un simple trait. Les $=X^{*}_{\ldots}$ du bas sont écrits sous les objets, avec un signe $=$ vertical.}

\marginal{\ill{} un \ill{} pas \ill{} vrai : \ill{} que $J\cap L=\emptyset$ on \ill{} \ill{} \ill{} \ill{} \ill{} \ill{} que $J\prec L$, \ill{} on posait \ldots\ $\Sigma_{J,L}=\dot T_J\cap\Sigma_L$ \ldots\ $\dot T_J\overset{?}{=}\varinjlim_{d\in\mathrm{Drap}(I)}\Sigma_d$ ($d$ \ill{} initial de $d\in J$) $(=\Sigma_{J,I})$ \ldots\ $\Sigma_{J,L}\overset{?}{=}\varinjlim\Sigma_d$ ($d\in\mathrm{Drap}(J\cup L)$, \ill{} initial de $d\in J$, \ill{} $\in L$)}\note{longues notes en travers dans toute la marge gauche, en partie biffées ; seules les formules se lisent à peu près, avec des points d'interrogation de sa main au-dessus des signes $=$.}

On a les inclusions
\[
\begin{array}{ccccc}
\Sigma_J & \subset & \Sigma_{J'} & \subset & \Sigma_{I'}\\
\| & & \| & & \|{}^{?}\\
X_J\smallsetminus X_J\cap\overset{\circ}{T}_{I'_0} & & X_{J'}\smallsetminus X_{J'}\cap\overset{\circ}{T}_{I'_0} & & X_{I'}\smallsetminus X_{I'}\cap\overset{\circ}{T}_{I'_0}
\end{array}
\]
\note{le point d'interrogation à côté du dernier signe $=$ est de sa main, mais sa lecture est incertaine.}

Considérons donc le voisinage tubulaire de $\Sigma_J$ dans $\Sigma_{I'}$

\page{86}
\note{En haut à gauche, « 40 » de sa main.}

\[
(75)\qquad T_{\Sigma_J,\Sigma_{I'}}=\Sigma_{I'}\cap\bigcup_{j\in J}T_j\quad\Big(=\Sigma_{I'}\cap\bigcup_{j\in\overline J}T_j\Big)
\]
\note{sous le second signe $=$, relié à lui par un trait : « mais en général c'est différent\ldots » ; le signe $=$ est donc mis en doute par lui.}

et considérons son bord
\[
(76)\qquad \dot T_{\Sigma_J,\Sigma_{I'}}=\varinjlim_{d\in\mathrm{Drap}(I')}(\Sigma_d)\qquad\big(\simeq\Sigma_{J,I'\smallsetminus J}\big)
\]
\note{sous la limite : « premier terme de $d$ est $\in J$, derniers termes sont dans $I'\smallsetminus J$ ». La parenthèse de droite ajoute « avec la notation ci-dessous (77) ». Un essai d'écriture en $\Sigma_{i_0,i_1,i_2}$ ($i_0\in J$, $i_1,i_2\in I'\smallsetminus J$) est barré de plusieurs traits à la suite de $(\Sigma_d)$.}

On pose (pour $L=I'\smallsetminus J$, \ill{} \ill{} \ill{})
\[
(77)\quad
\begin{cases}
\Sigma_{J,L}=\dot T_{\Sigma_J,\Sigma_{I'}}\smallsetminus\underbrace{\dot T_{\Sigma_J,\Sigma_{I'}}\cap\underbrace{\overset{\circ}{T}_{J'}}_{\bigcup_{j\in J'}\overset{\circ}{T}_j}}_{\dot T_{\Sigma_J,\Sigma_{I'}}\cap\bigcup_{j\in J'}\overset{\circ}{T}_j}\\[1ex]
\qquad=\dot T_{\Sigma_J,\Sigma_{I'}}\cap\Sigma_L
\end{cases}
\]
\note{ici et dans (78)-(79), $\Sigma_{J,L}$ porte un petit exposant biffé, qui n'a pas pu être lu. Dans $\overset{\circ}{T}_{J'}$ et sous la réunion, une barre sur $J'$ est biffée.}

\marginal{NB $\Sigma_L=\Sigma_{I'}\smallsetminus\Sigma_{I'}\cap\overset{\circ}{T}_{J'}$}\note{en travers dans la marge gauche, en face de (77).}

On aura sauf erreur
\[
(78)\qquad \Sigma_{J,L}\simeq\varinjlim_{d\in\mathrm{Drap}(J\cup L)}\Sigma_d
\]
\note{sous la limite : « premier terme de $d$ dans $J$, derniers termes dans $L$ ».}

On a
\[
(79)\qquad
\begin{array}{ccc}
\Sigma_{J,L} & \hookrightarrow & \Sigma_L\\
\downarrow & & \\
\Sigma_J & &
\end{array}
\]
\note{sous la flèche horizontale : « inclusion canonique » ; à gauche de la flèche verticale : « induite par $\dot T_{\Sigma_J\Sigma_{I'}}\to\Sigma_J$ ».}

\page{87}
\note{Sans numéro de sa main ; suite du feuillet 40.}

qui s'explicitent en termes de (78) en notant que tout drapeau de $J\cup L$ \add{intervenant dans (78)} s'écrit de façon unique
\[
(80)\qquad (j_1\ldots j_p\,l_1\ldots l_q)\qquad j_1\ldots j_p\in J,\quad l_1\ldots l_q\in L,\quad p,q\ge 1
\]
et \struck{\ill{}} on a des
\[
(81)\qquad
\begin{array}{ccc}
\Sigma_{j_1\ldots j_p l_1\ldots l_q} & \hookrightarrow & \Sigma_{l_1\ldots l_q}\\
\downarrow & & \\
\Sigma_{j_1\ldots j_p} & &
\end{array}
\]
\note{la page 87 passe de (79) à (80) et (81) sans lacune.}

\emph{Remarque}. Les $\Sigma_{J,L}$ ont un sens \add{(dans une catégorie isotopique \ill{})} indépendamment de rétractions coniques locales autour des $X_i^{*}$, du moins en théorie topologique modérée, ainsi que ces $\Sigma_{J,L}\hookrightarrow\Sigma_L$ -- et la description que j'en ai donnée ici est adaptée à ce cas -- la structure $\Sigma_{J,L}\to\Sigma_J$ par contre utilise les \add{rétractions} \struck{structures côniques} locales. Quand on dispose de structures côniques locales (ce qui, en topologie modérée, et dans le cas équisingulier, est une donnée supplémentaire \ill{} « \ill{} »), il est plus commode,\note{la parenthèse « (dans une catégorie isotopique \ldots) » est ajoutée au-dessus de la ligne et entourée ; au-dessus de « côniques » biffé, deux mots, biffés aussi, n'ont pas été lus. Il écrit ici « rétractions coniques » sans accent, mais « structures côniques » avec. La page s'arrête sur « plus commode, » ; la suite ne se trouve pas à la page 88, qui ouvre le paragraphe 17.}

\page{88}
\note{En haut à gauche, « 41 » de sa main.}

\subsection*{17. Récapitulation sur les $\mathcal V_i$, $\mathcal V_{i_0\ldots i_n}$, $\mathcal V_J$, $\mathcal V_J^L$, $\mathcal V_{J,L}$, et introduction des $\mathcal V_{J_0,J_1,\ldots J_n}$}

Soit $I$ un ens.\ ordonné. Une partie \add{$J$} de $I$ est dite \emph{fermée} si $j\in J$, $i\le j\Rightarrow i\in J$. Les parties fermées de $I$ forment une \struck{ens.} partie de $\mathfrak P(I)$ stable par \struck{intersection,} réunions \struck{finies,} quelconques, et par \struck{intersections} \add{intersections} quelconques, ce sont donc les fermés d'une topologie sur $I$, appelée la \emph{topologie} \struck{\ill{} \ill{}} \emph{inférieure}. \struck{Soit $L=I\smallsetminus J$, \ill{}} On a donc pour une partie $A$ de $I$
\[
(82)\qquad \overline A=\lbrace i\in I\mid \exists j\in A \text{ tel que } i\le j\rbrace
\]
(adhérence de $A$ pour la top.\ inférieure de $I$)\note{dans la phrase « par \ldots\ intersections quelconques », le premier « intersection » biffé est en fin de ligne, le mot est récrit en surcharge ; « inférieure » est écrit au-dessus d'un mot biffé qui semble « coarctique » ou analogue, non lu.}

\struck{Si} Soit $A\subset I$, $L=I\smallsetminus A$, pour que $J$ soit fermé, \add{i.e.\ $L$ soit \emph{ouvert},} il faut et il suffit que $L$ soit « \emph{cofermé} » i.e.\ fermé pour la top.\ associée à l'ordre opposé à $I$ \add{celui} de $I$, i.e.\ $l\in L$, \struck{$i\in\ldots$} $i\ge l\Rightarrow i\in L$. L'adhérence d'une partie $B$ de $I$ pour la top.\ supérieure est notée $\widetilde B$ :
\[
(83)\qquad \widetilde B=\lbrace i\in I\mid \exists j\in B \text{ tel que } i\ge j\rbrace
\]
\struck{Donc} Ainsi $\widetilde B$ est aussi le plus petit ouvert (pour la top.\ inf.)\ contenant $B$.\note{« pour que $J$ soit fermé » : il faut sans doute lire $A$ ; la lettre est récrite. Dans « opposé à $I$ \add{celui} de $I$ », la lecture est incertaine.}

\struck{Une} Pour une partie $J$ de $I$, il revient au même qu'elle soit loc.\ fermée (i.e.\ intersection d'un fermé et d'un ouvert), ou qu'on ait $J=\overline J\cap\widetilde J$, ou encore

\page{89}
\note{Sans numéro de sa main ; suite du feuillet 41.}

que \struck{\ill{}} $(j,k\in J$ et $j\le i\le k)\Rightarrow i\in J$. On dit \add{(\ill{} $j\le i\le k$)} aussi, dans ce cas, que $J$ est une \emph{bande} de $I$.

\emph{C'est} \struck{\ill{} \ill{}} \add{\ill{}} \ill{} par passage de $I$ à $I^{\circ}$. Toute intersection de bandes est une bande. La plus petite bande contenant $J$ est \struck{$\overline J\cap\widetilde J$} \ill{}. Soient $J$, $K$ deux \struck{bandes} parties de $I$, on écrit
\[
(85)\qquad J\prec K\ \overset{\text{déf}}{\Longleftrightarrow}\ J\subset\overline K,\ K\subset\widetilde J
\]
(relation de \emph{domination}). C'est une relation de préordre sur l'ens.\ $\mathfrak P(I)$ des parties de $I$, la relation d'équivalence associée est
\[
(86)\qquad J\sim K\iff J\subset\underbrace{\overline K\cap\widetilde K}_{B(K)},\ K\subset\underbrace{\overline J\cap\widetilde J}_{B(J)}\iff B(J)=B(K)
\]
L'ens.\ ordonné quotient est l'ens.\ des \emph{bandes} de $I$, ordonné par la relation de \struck{\ill{}} domination.

\marginal{(84) $\overline J\cap\widetilde J=\lbrace i\in I\mid \exists j,k\in J,\ j\le i\le k\rbrace\overset{\text{déf}}{=}B(J)$}\note{en travers dans la marge gauche, en face des lignes « C'est \ldots\ La plus petite bande contenant $J$ est \ldots » ; une accolade la rattache à cette phrase. C'est le seul lieu où figure le numéro (84).}

Supposons maintenant que $J$ et $K$ soient des bandes, \add{avec $J\prec K$,} soit $I'$ la bande engendrée par $J\cup K$, on a
\[
(87)\qquad I'=\lbrace i\in I\mid \exists j\in J,\ k\in K \text{ t.q. } j\le i\le k\rbrace=\widetilde J\cap\overline K
\]
et $J$ est fermé dans $I'$, $K$ est ouvert dans $I'$, \struck{et} en fait $K$ est dense dans $I'$ et $J$ codense dans $I'$ i.e.\ $I'=\overline K$ et $I'\subset\widetilde J$ \struck{\ill{}} i.e.\ $I'\subset\overline K\cap\widetilde J$ (on a \ill{} égalité\,!). Inversement si $\exists$ bande $I'$ telle que $J$ soit un fermé codense dans $I'$, $K$ un ouvert dense dans $I'$, alors $J$ et $K$ sont des bandes, $J\prec K$, et on a $I'=\widetilde J\cap\overline K$.\note{« $I'=\overline K$ » : il faut sans doute lire $I'\subset\overline K$ (le signe est écrit gras) ; lu ici tel quel.}

\page{90}
\note{En haut à gauche, « 42 » de sa main.}

\struck{Une filtration \add{\struck{ordinaire}} de type $I$} Soient $X$ un espace, $I$ un ens.\ ordonné. Une \emph{stratification ordinaire}, de type $I$, de $X$, est une application
\[
I\longrightarrow\mathfrak P(X)\qquad i\longmapsto X_i
\]
satisfaisant les \ill{} conditions
\begin{itemize}
\item[(Strat 1)] les $X_i$ fermés, la famille des $X_i$ \add{est un \emph{recouvrement} loc.\ fini de $X$} ;
\item[(Strat 2)] $i\le j\Rightarrow X_i\subset X_j$ \quad i.e.\ $i\mapsto X_i$ est croissante ;
\item[(Strat 3)] $\forall i,j\in I$, on a \quad $X_i\cap X_j=\bigcup_{k\in\bar i\cap\bar j}X_k$.
\end{itemize}
\note{« les deux conditions » : il en énumère trois. En travers dans la marge gauche, en face des conditions : « les $X_i$ \ill{} \ill{} les \emph{fermés} » (lecture incertaine).}

\struck{Posons les conditions} Pour tout drapeau
\[
d=(i_0<\cdots<i_n)
\]
de $I$, posons
\[
(88)\qquad X_d=X_{i_0},
\]
alors \struck{$\varinjlim_{d\in\mathrm{Drap}\,I}$} $X_d$ forment un \struck{syst.} espace semi-simplicial \struck{\ill{}} (pour des applications strictement croissantes sur les simplexes types) \add{au-dessus de $X$}, qui forme en un sens convenable une « résolution » semi-simpl.\ de $X$. \uncertain{Ainsi} p.\ ex.\ pour tout faisceau abélien $F$ sur $X$
\[
(89)\qquad H^{*}(X,F)\Longleftarrow E_2^{pq}=H^p\big(I,\ i\mapsto H^q(X_i,F)\big)=H^p_{ss}\big(\mathrm{Drap}(I),\ d\mapsto H^q(X_d,F)\big)
\]
Mais les \add{$X_i$,} $X_d$ sont des espaces souvent compliqués, et la dévissage (89) souvent assez peu pratique, en remplaçant par les $X_i^{*}$, $X_d^{*}=X^{*}_{i_0}$, on pourrait

\page{91}
\note{Sans numéro de sa main ; suite du feuillet 42.}

\[
(90)\qquad X_i^{*}=X_i\smallsetminus\bigcup_{j<i}X_j
\]
Les $X_i^{*}$ non \emph{vides} forment alors une partition de $X$. Les $X_i^{*}$ sont les « strates ouvertes ».\note{devant $X_i$, dans (90), un signe biffé.}

On introduit parfois des axiomes supplémentaires :
\begin{itemize}
\item[(Strat 4)] (fidélité) \quad $\forall i\in I$, $X_i^{*}\ne\emptyset$.
\end{itemize}
Cet axiome a l'inconvénient de ne pas être stable par restriction à un ouvert de $X$. Il est parfois complété par l'axiome plus fort encore moins stable
\begin{itemize}
\item[(Strat 4 bis)] (Canonicité de $I$) les $X_i^{*}$ sont connexes non vides, i.e.
\item[(Strat 5)] (densité) $\forall i\in I$, $X_i^{*}$ dense dans $X_i$, \add{i.e.\ $X_i=\overline{X_i^{*}}$,} i.e.\ $j<i\Rightarrow X_j$ \emph{rare} dans $X_i$.
\end{itemize}
\note{la ligne de (Strat 4 bis) s'achève sur « i.e. » et la ligne suivante porte (Strat 5) : on ne sait si (Strat 5) est l'explicitation de (Strat 4 bis) ou un axiome distinct. « (densité) $\forall i\in I$ » et « i.e.\ $X_i=\overline{X_i^{*}}$ » sont ajoutés au-dessus de la ligne.}

\marginal{moyennant (Strat 5), la condition de fidélité équivaut : $\forall i\in I$, $X_i\ne\emptyset$ (car $X_i=\overline{X_i^{*}}$ donc $X_i\ne\emptyset\Rightarrow X_i^{*}\ne\emptyset$)}\note{en travers dans la marge gauche, en face de (Strat 5).}

Cet axiome a l'avantage d'être stable par restriction à un ouvert. Mais on ne supposera (Strat 4) (voire Strat 4 bis) ou (Strat 5) que par mention expresse.

\struck{Stratifications de fidélité \ill{}}

On a pour tout $i\in I$
\[
(91)\qquad X_i=\bigcup_{j\in\bar i}X_j^{*}\qquad(\text{i.e. } j\le i)
\]
donc \uncertain{on aura}, si $i,j\in J$\note{la suite de cette page est la page 93 (sans numéro, avec les formules (92) et (93)), et non la page 92 (feuillet 43, qui s'ouvre sur (94)) : les deux feuillets ont été numérisés dans l'ordre 92, 93, mais se lisent 91, 93, 92.}

\page{92}
\note{En haut à gauche, « 43 » de sa main.}

\struck{$X^{*}_J=X_{\overline J}$}

\note{ligne barrée en tête de page. Ce feuillet fait suite à la page 93, où $X^{*}_J$ est défini par (93).}

\uncertain{Les} \ill{} compatibilité de $J\mapsto X^{*}_J$ \ill{} \ill{} réunions et intersections
\[
(94)\qquad \bigcup_\alpha X^{*}_{J_\alpha}=X^{*}_{\bigcup_\alpha J_\alpha},\qquad \bigcap_\alpha X^{*}_{J_\alpha}=X^{*}_{\bigcap_\alpha J_\alpha}
\]
en particulier
\[
X^{*}_\emptyset=\emptyset,\qquad X^{*}_I=X.
\]
On a
\[
(95)\qquad X^{*}_J\subset X_J=X^{*}_{\overline J}
\]
\note{un $X$ biffé devant $X^{*}_{\overline J}$ ; le signe d'inclusion est surchargé et ressemble à $\in$.}

\marginal{en particulier si $J$ fermé, on a $X_J=X^{*}_J$, et $X^{*}_J$ est fermé}\note{à droite de (95).}

et
\[
(96)\qquad X^{*}_J=X^{*}_{\overline J}\smallsetminus X^{*}_{\underbrace{\scriptstyle\overline J\smallsetminus J}_{J_0}}
\]
Si $J$ est une bande, $J_0$ est fermé donc $X^{*}_J$ est loc.\ fermé.

\ill{} \ill{} (Strat 5), \uncertain{alors}
\[
(97)\qquad \overline{X^{*}_J}=X_J\ (=X^{*}_{\overline J})=X^{*}_J\amalg X^{*}_{\overline J\smallsetminus J}
\]
\note{à droite de (97) : « moyennant \emph{densité} ».}

\struck{$\dot X^{*}_J\overset{\text{déf}}{=}\overline{X^{*}_J}\smallsetminus X^{*}_J$}

\marginal{dans \ill{} $X^{*}_J$ \ill{} le plus petit \ill{} \ill{} de $X$ \ill{} de la forme $X^{*}_?$, \ill{} \ill{} soit \ill{} (97) \ill{} $X^{*}_J$ \ill{} ouvert \ill{} \ill{}}\note{en travers dans la marge gauche, en face de (97) ; lecture très incertaine.}

donc \struck{$\overline{X^{*}_J}$ \ill{}} $X^{*}_J$ est fermé i.e.\ égal à $\overline{X^{*}_J}$ ($=X^{*}_{\overline J}$) ssi $X^{*}_{\overline J\smallsetminus J}=\emptyset$ i.e.\ si les $X_i^{*}$ ($i\in\overline J\smallsetminus J$) sont vides, i.e.\ (si strat.\ est fidèle) \emph{ssi} $\overline J\smallsetminus J=\emptyset$ i.e.\ $J=\overline J$ i.e.\ $J$ fermé.

Si $L\subset I$ est ouvert, \add{i.e.\ $J=I\smallsetminus L$ fermé,} alors
\[
(98)\qquad X^{*}_L=\complement_X X^{*}_J
\]
est ouvert (si $L\subset I$ est ouvert)\note{« ssi » est doublement souligné. Au début de la ligne « Si $L\subset I$ est ouvert », une lettre biffée.}

\page{93}
\note{Sans numéro de sa main ; suite de la page 91 (voir la note en fin de page 91) : cette page précède, dans la lecture, le feuillet 43 (page 92).}

\[
X_i\subset X_j\iff\forall i'\in\bar i,\ X^{*}_{i'}\subset X_j=\bigcup_{i''\in\bar j}X^{*}_{i''}
\]
$\iff$ $\forall i'\in\bar i$ t.q.\ $X^{*}_{i'}\ne\emptyset$, on a $i'\in\bar j$ i.e.\ $i'\le j$.
Si la stratification est fidèle, alors on a donc
\[
(92)\qquad X_i\subset X_j\iff X_i^{*}\subset X_j\iff i\le j
\]
(si strat.\ fidèle)\note{la première ligne se lit sur deux lignes, la seconde équivalence sous la première ; les indices sont surchargés.}

ce qui implique que $I\to\mathfrak P(X)$ est injective, et que la relation d'ordre de $I$ est induite par la relation d'inclusion de $\mathfrak P(X)$. (\ill{} \ill{} peut identifier $I$ à une partie de $\mathfrak P(X)$), et la donnée d'une str.\ fidèle \emph{équivaut} à la donnée d'une \struck{\ill{}} famille \add{loc.\ finie} de parties fermées \add{non vides} de $X$, \add{stable par l'intersection,} satisfaisant la condition (Strat 3) \add{les « strates fermées » ; de réunion $X$} et la condition qu'aucune strate fermée n'est réunion de ses sous-strates strictement plus petites. Si on a Strat 5, cet \struck{\ill{}} axiome \ill{} plus fort que \ill{} les strates fermées.\note{les ajouts interlinéaires de ce paragraphe sont reliés au texte par des traits ; leur place exacte est incertaine, de même que la fin du paragraphe.}

Si $J$ est une partie de $I$, on pose
\[
(93)\qquad
\begin{cases}
X^{*}_J=\bigcup_{j\in J}X^{*}_j & (\text{pas nécessairement fermé})\\
X_J=\bigcup_{j\in J}X_j & \text{fermé dans } X
\end{cases}
\]
\struck{Nous utiliserons \ill{} les notations suivantes si $J$ est une bande. Alors on a \ill{}}

\marginal{NB $J\mapsto X^{*}_J$ est compatible aux réunions, \ill{} intersections : $X^{*}_{J\cap K}=X^{*}_J\cap X^{*}_K$, $X^{*}_{J\cup K}=X^{*}_J\cup X^{*}_K$}\note{note encadrée, en travers dans le coin inférieur gauche, reliée par une flèche à (93) ; lecture incertaine dans le détail.}

\page{94}
\note{En haut à gauche, « 44 » de sa main : suite du feuillet 43 (page 92), qui s'arrête sur (98). Les numéros (94) à (97) de cette page répètent ceux des pages 92-93 pour d'autres formules ; ils sont gardés tels qu'écrits.}

et moyennant densité et fidélité, la \struck{\ill{}} \add{\uncertain{réciproque}} est vraie : $X^{*}_L$ est ouvert \emph{ssi} $L$ est ouvert. Sous les mêmes hypothèses, pour que $X^{*}_J$ soit loc.\ fermé i.e.\ $\overline{X^{*}_J}\smallsetminus X^{*}_J$ \add{$=X^{*}_{\overline J\smallsetminus J}$} soit fermé, \struck{i.e.} il faut et il suffit que $\overline J\smallsetminus J$ soit fermé i.e.\ que $J$ soit loc\textsuperscript{t} fermé i.e.\ soit une bande.

\struck{Soient les}

Soit maintenant, pour tout $i\in I$, $\mathcal V_i$ le pro-espace
\[
(94)\qquad \mathcal V_i=\mathcal V(X_i^{*},X)
\]
\note{à droite : « donc $X_i^{*}\overset{(h)}{\hookrightarrow}\mathcal V_i$ équivalence d'homotopie ($X$ unibranche \ill{}) » ; lecture incertaine de la parenthèse.}

et considérons $\mathcal V_i\cap X_j^{*}$, \ill{}.
\[
(95)\qquad \mathcal V_i\cap X_j^{*}\ne\emptyset\iff\mathcal V_i\cap\overline{X_j^{*}}\ne\emptyset\iff X_i^{*}\cap\overline{X_j^{*}}\ne\emptyset
\ \overset{(a)}{\Longrightarrow}\ X_i^{*}\cap X_j\ne\emptyset\ \overset{(b)}{\Longrightarrow}\ i\le j
\]
\note{sous la flèche (a) : « équiv.\ si strat.\ dense » ; sous la flèche (b) : « équiv.\ si strat.\ fid. » (les lettres (a), (b) sont de l'éditeur). Sous le numéro (95), un signe biffé.}

donc dans le cas d'une strat.\ dense et fidèle,
\[
\mathcal V_i\cap X_j^{*}\ne\emptyset\iff i\le j.
\]
\struck{\ill{} \ill{}} \ill{}.

On fera attention que cette relation n'est pas \ill{} \ill{} symétrique en $(i,j)$, \struck{\ill{}} donc n'est pas de \ill{} \ill{} d'être « équivalente » : $\mathcal V_i\cap\mathcal V_j\ne\emptyset$. On a
\[
(96)\qquad \mathcal V_i\cap\mathcal V_j\ne\emptyset\ \overset{(c)}{\Longrightarrow}\ i \text{ et } j \text{ comparables i.e. } i\le j \text{ ou } j\le i
\]
et la réciproque ($\Longleftarrow$) si strat.\ fidèle.\note{au-dessus de la flèche (c) : « si $X$ métrisable, dans \ill{} \ill{} \ill{} \ill{} aux $X$ \ill{} » ; lecture très incertaine. La flèche réciproque, sous la première, porte « strat.\ fidèle ». Le numéro, écrit (95) ou (96), se lit mal.}

On posera
\[
(97)\qquad \mathcal V_{ij}\overset{\text{déf}}{=}\mathcal V_i\cap\mathcal V_j,
\]

\page{95}
\note{Sans numéro de sa main ; suite du feuillet 44.}

en faisant attention que pour $i$ et $j$ incomparables, si $X$ unibranche, on trouve \struck{\ill{}} donc $\mathcal V_{ij}=\emptyset$ \ill{} (dans \uncertain{par là} \uncertain{fixés} d'un \uncertain{coup}\,!), et pour $i=j$, on trouve $\mathcal V_{ij}=\mathcal V_{ii}=\mathcal V_i$ (dans \uncertain{par là} \uncertain{besoin} de l'introduire, $\mathcal V_i$ suffit\,!), enfin \struck{\ill{}} si $i$ et $j$ sont distincts et comparables, \struck{\ill{}} comme $\mathcal V_{ij}=\mathcal V_{ji}$, il suffit d'introduire $\mathcal V_{ij}$ pour \struck{\ill{}} $i<j$. Plus généralement, on introduit les
\[
(98)\qquad \mathcal V_{i_0\ldots i_n}=\mathcal V_{i_0}\cap\cdots\cap\mathcal V_{i_n},
\]
mais on se bornera en règle générale au cas où $\lbrace i_0,\ldots,i_n\rbrace$ est un drapeau (sinon on trouve $\emptyset$) et où on fait \uncertain{les indices} à \uncertain{jamais} en ordre str.\ croissant
\[
i_0<i_1<\cdots<i_n
\]
pour éviter des redondances). Les pro-espaces${}^{*}$ sont les « pièces de construction » élémentaires, pour une construction des \uncertain{types} d'homotopie, par $\varinjlim$ (finies si $I$ fini), compte tenu des morphismes semi-simpliciaux évidents entre les $\mathcal V_d$ ($d=(i_0,\ldots,i_n)\in\mathrm{Drap}(I)$).\note{(98) répète le numéro de la page 92 pour une autre formule. La phrase qui entoure les parenthèses se lit mal ; plusieurs mots n'ont été lus qu'avec doute.}

\marginal{${}^{*}$ (et de \uncertain{tels} pro-\uncertain{espaces} $\mathcal V_{i_0,\ldots,i_n}$ \ill{} \ill{} \ill{})}\note{note de renvoi en travers dans la marge gauche, rattachée par l'astérisque à « pro-espaces » ; lecture incertaine.}

\page{96}
\note{En haut à gauche, « 45 » de sa main.}

On pose, si $Z$ est une partie de $X$,
\[
(99)\qquad \mathcal V^Z_{i_0,\ldots,i_n}=\mathcal V_{i_0\ldots i_n}\cap Z\qquad(\hookrightarrow\mathcal V_{i_0\ldots i_n})
\]
si $Z$ est \struck{fermé} une partie loc\textsuperscript{t} fermée qui contient \struck{\ill{}} $X^{*}_{i_n}$ \struck{$X_{i_n}\smallsetminus\dot X_{i_0}$ (NB $\dot X_{i_0}\overset{\text{déf}}{=}X_{i_0}\smallsetminus X^{*}_{i_0}$ $=\bigcup_{i<i_0}X_i$ $=\bigcup_{i<i_0}X^{*}_i$)} \add{(de sorte que $X_{i_n}\overset{(h)}{\hookrightarrow}\mathcal V^Z_{i_n}$ $\overset{(h)}{\hookrightarrow}\mathcal V_{i_n}$)}, alors on a
\[
(100)\qquad \mathcal V^Z_{i_0\ldots i_n}\overset{(h)}{\hookrightarrow}\mathcal V_{i_0\ldots i_n}\qquad(\text{équiv.\ d'homotopie})
\]
(on suppose $X$ unibranche\ldots)\note{la ligne barrée l'est d'un long trait. Au-dessus de « alors », un point d'interrogation. L'ajout entre parenthèses est écrit au-dessus de la ligne barrée ; son premier terme se lit mal ($X_{i_n}$ ou $X^{*}_{i_n}$).}

\marginal{? à vérifier\,!}\note{dans la marge gauche, en face de (100) ; « vérifier » est souligné.}

On va généraliser la définition des $\mathcal V_i$, $\mathcal V_{i_0\ldots i_n}$ en des $\mathcal V_J$, $\mathcal V_{J_0\ldots J_n}$, où $J$ est une bande de $I$, resp.\ $J_0\prec\cdots\prec J_n$ un « drapeau de bandes » de $I$. On définit
\[
(101)\qquad \mathcal V_J=\mathcal V(X^{*}_J,X)=\bigcup_{j\in J}\mathcal V_j=\varinjlim_{d\in\mathrm{Drap}(J)}\mathcal V_d
\]
\note{à droite de la première égalité : « (pro-voisinage tubulaire) » (lecture incertaine du premier mot).}

(NB la définition et les deux formules \uncertain{marchent} \uncertain{encore} pour $J$ partie quelconque de $I$)
\[
(102)\qquad \mathcal V^Z_J=\mathcal V_J\cap Z
\]
On a donc
\[
(103)\qquad X^{*}_J\overset{(h)}{\hookrightarrow}\mathcal V_J\qquad(\text{équiv.\ d'hom.})
\]
\[
(104)\qquad X^{*}_J\overset{(h)}{\hookrightarrow}\underbrace{\mathcal V^Z_J}_{=\mathcal V(X^{*}_J,Z)}\overset{(h)}{\hookrightarrow}\mathcal V_J\qquad(\text{équiv.\ d'hom.})\ \text{si } Z\supset X_J
\]
\note{« si » est doublement souligné ; sous $\mathcal V^Z_J$, un signe $=$ vertical et « $\mathcal V(X^{*}_J,Z)$ ».}

\page{97}
\note{Sans numéro de sa main ; suite du feuillet 45.}

On aura
\[
(105)\qquad \mathcal V^Z_J=\bigcup_{j\in J}\mathcal V^Z_j=\varinjlim_{d\in\mathrm{Drap}(J)}\mathcal V^Z_d
\]
\note{sous $\mathcal V^Z_d$, un signe $\simeq$ vertical marqué (h) le relie à $\mathcal V_d$, avec la condition « si $Z\supset X^{*}_J$ ». Plus bas, un mot biffé.}

On pose, pour deux bandes $J$ et $L$
\[
(106)\quad
\begin{cases}
\mathcal V_{J,L}=\mathcal V_J\cap\mathcal V_L, & \mathcal V^Z_{J,L}=\mathcal V_{J,L}\cap Z\\
\mathcal V^{L}_{J}=\mathcal V_J\cap X^{*}_L=\mathcal V_{J,L}\cap X^{*}_L=\mathcal V^{X^{*}_L}_{J,L}
\end{cases}
\]
\note{la seconde ligne : un exposant biffé après le $L$ de $\mathcal V^L_J$ ; sous $\mathcal V_J\cap X^{*}_L$, relié par un signe $=$ vertical, « $\mathcal V^{X^{*}_L}_J$ \ldots\ notation (102) » ; sous le dernier terme, « avec $Z=X^{*}_L$ », dont l'exposant $X^{*}_L$ est récrit sur un $Z$.}

\marginal{Suggestion : $\mathcal V^{\flat}_{J,L}$ au lieu de $\mathcal V^L_J$ \ldots\ $\mathcal V^{\flat}_{i_0\ldots i_n}=\mathcal V^{X^{*}_{i_n}}_{i_0\ldots i_n}=\mathcal V_{i_0}\cap\ldots\cap\mathcal V_{i_{n-1}}\cap X^{*}_{i_n}$}\note{dans la marge gauche, en face de (106), reliée à l'accolade par un trait. Le signe en exposant ressemble à un bémol ; lecture incertaine.}

On aura donc
\[
(107)\qquad \mathcal V_{J,L}=\bigcup_{\substack{j\in J\\ l\in L}}\underbrace{\mathcal V_j\cap\mathcal V_l}_{\mathcal V_{j,l}}
=\bigcup_{i\in J\cap L}\mathcal V_i\ \cup\bigcup_{\substack{j<l\\ j\in J,\ l\in L}}\mathcal V_{jl}\ \cup\bigcup_{\substack{l<j\\ j\in J,\ l\in L}}\mathcal V_{lj}
\]
\note{après le premier membre, une réunion biffée, indexée par « $j\in J$, $l\in L$, $j\le l$ ».}

\marginal{$=\varinjlim_{\substack{d\in\mathrm{Drap}(J\cup L)\\ d\cap J\ne\emptyset,\ d\cap L\ne\emptyset}}\mathcal V_d$}\note{dans la marge gauche, en face de la seconde ligne de (107).}

\struck{$\varinjlim_{d\in\mathrm{Drap}_{\ge 2}(J\cup L)}\mathcal V_d$} donc si $J\prec L$, et $J$ et $L$ des bandes
\[
(108)\qquad \mathcal V_{J,L}=\bigcup_{i\in J\cap L}\mathcal V_i\ \cup\bigcup_{\substack{j<l\\ j\in J,\ l\in L}}\mathcal V_{jl}
=\varinjlim_{\substack{d\in\mathrm{Drap}(J\cup L)\\ d\cap J\ne\emptyset,\ d\cap L\ne\emptyset}}\mathcal V_d
\]
\note{la limite barrée l'est de plusieurs traits obliques ; son indice se lit mal.}

\marginal{NB ici les drapeaux qui interviennent $i_0<\cdots<i_n$, tous $i_\alpha\in J\cup L$, $i_0\in J$, $i_n\in L$, donc les $i_\alpha$ qui sont $\in J$ forment une \uncertain{segment} initial, les $i_\beta\in L$ un segment final${}^{*}$}\note{en bas à droite, sous (108).}

\marginal{${}^{*}$ la réunion des deux segments \ldots\ $\lbrace i_0,\ldots,i_n\rbrace$ ; l'intersection \ill{} $\emptyset$ si $J\cap L=\emptyset$\ldots}\note{en bas à gauche, rattachée à la note précédente par l'astérisque ; lecture incertaine.}

\page{98}
\note{En haut à gauche, « 46 » de sa main. Dernière page du dossier.}

On aura de même, pour $J$ et $L$ quelconques
\[
(109)\qquad \mathcal V^Z_{J,L}=\varinjlim_{\substack{d\in\mathrm{Drap}(J\cup L)\\ d\cap J\ne\emptyset,\ d\cap L\ne\emptyset}}\mathcal V^Z_d
\]
\note{sous $\mathcal V^Z_d$, une flèche verticale marquée $\simeq$ (h) vers $\mathcal V_d$, avec la condition : « si $J\prec L$, $J$ et $L$ des bandes et si $Z$ contient $X^{*}_L$, donc contient $X^{*}_{i_n}$ si $d=(i_0,\ldots,i_n)$ » (un mot biffé devant $X^{*}_L$).}

donc si \add{$J,L$ des bandes et $J\prec L$, et si} $Z$ contient \struck{\ill{}} $X^{*}_L$, \ill{},
\[
(110)\qquad \mathcal V^Z_{J,L}\overset{(h)}{\hookrightarrow}\mathcal V_{J,L}
\]
\note{devant « donc si », un mot biffé. Le $J$ du second membre est surchargé.}

en particulier
\[
(111)\qquad \mathcal V^L_J\ \Big(\overset{\text{déf}}{=}\mathcal V^{X^{*}_L}_{J,L}\Big)\overset{(h)}{\hookrightarrow}\mathcal V_{J,L}
\]
Soient
\[
(112)\qquad J_0\prec J_1\prec\cdots\prec J_n
\]
des bandes, on pose
\[
(113)\qquad \mathcal V_{J_0,\ldots,J_n}=\mathcal V_{J_0}\cap\cdots\cap\mathcal V_{J_n}
=\bigcup_{\substack{j_0\in J_0\\ \cdots\\ j_n\in J_n}}\mathcal V_{j_0 j_1\ldots j_n}
=\varinjlim_{\substack{d\in\mathrm{Drap}(J_0\cup J_1\cup\cdots\cup J_n)\\ d\cap J_0\ne\emptyset,\ldots,\ d\cap J_n\ne\emptyset}}\mathcal V_d
\]

\struck{\ill{}}\note{au bas de la page, un bloc de formules, raturé en tous sens, n'a pas pu être lu. Le dossier s'arrête ici.}

\end{document}
